\documentclass{article}
\usepackage{mathnotes}

\title{Binomial Theorem}
\description{Binomial Theorem and Coefficient}

\begin{document}

\begin{definition}[Binomial Coefficient]
For $n, k \in \naturals, n \geq k \geq 0,$ the \textbf{binomial coefficient} is the coefficient of the $x^k$ term in the polynomial expansion of the binomial power $(1+x)^n.$ It is denoted as $\binom{n}{k}$ and can be expressed through \@{factorial} notation as

\[ \binom{n}{k} = \frac{n!}{k!(n - k)!}. \]
\end{definition}

\begin{definition}[Generalized Binomial Coefficient]
For any complex number $n \in \complexes$ and any \@{integer} $k \geq 0,$ the \textbf{generalized binomial coefficient} is defined by

\[ \binom{n}{k} = \frac{n(n-1)(n-2) \cdots (n - k + 1)}{k!}. \]
\end{definition}

\begin{theorem}[Binomial Theorem]\label{binomial-theorem}
The expansion of any nonnegative integer power $n$ of the \@{binomial} $x+y$ is a sum of the form

\[ (x+y)^n = \sum_{k=0}^n \binom{n}{k} x^{n-k} y^k = \sum_{k=0}^n \binom{n}{k} x^k y^{n-k}. \]
\end{theorem}

\detach

\begin{remark}
A special case of the \@{binomial-theorem} is the case where $y = 1:$

\[ (x+1)^n = \sum_{k=0}^n \binom{n}{k} x^k. \]
\end{remark}

\end{document}
